% TOP_PROBLEM: {"id": 10000042, "title": "Percolation nonuniqueness on products with the line", "category_id": 19, "category": {"id": 19, "name": "probability", "display_name": "Probability", "description": "Problems involving probability theory, stochastic processes, and random structures.", "slug": "probability", "order_index": 19, "created_at": "2026-07-31T15:26:25.671Z"}, "status": "open", "problem_number": "AMR-099-0042", "set_id": 13}
% FIRST_POSTED: 2026-10-01
% ENTRY_KIND: statement_only
% UnsolvedMath ID 10000042; source catalogue code AMR-099-0042
% !TeX program = lualatex
% Definition and sources only; this file is not an LLM solution attempt.
\documentclass[11pt,a4paper]{article}
\usepackage[margin=25mm]{geometry}
\usepackage{fontspec}
\setmainfont{lmroman10-regular.otf}[BoldFont=lmroman10-bold.otf,
 ItalicFont=lmroman10-italic.otf,BoldItalicFont=lmroman10-bolditalic.otf]
\usepackage{amsmath,amssymb,mathtools}
\usepackage{unicode-math}
\setmathfont{latinmodern-math.otf}
\AtBeginDocument{\let\setminus\smallsetminus}
\usepackage{xurl,hyperref}
\hypersetup{colorlinks=true,urlcolor=blue}
\setcounter{secnumdepth}{0}
\setlength{\parindent}{0pt}
\setlength{\parskip}{5pt}
\setlength{\emergencystretch}{3em}
\begin{document}
\section*{Percolation nonuniqueness on products with the line}\label{10000042}
{\small UnsolvedMath ID 10000042 · AMR-099-0042 · Probability · AMR Open Problem Lists}
\subsection{Definitions and mathematical statement}
For an infinite graph $H$, define its Cheeger constant by $h(H)=\inf_{0<|A|<\infty}|\partial_H A|/|A|$, using external vertex boundary. Call a connected locally finite graph $G$ strongly amenable in this problem's sense if no infinite connected subgraph $H\subseteq G$ has $h(H)>0$. This is a graph-theoretic definition, not a condition on group actions.

The Cartesian product $G\times\mathbb Z$ has vertices $(v,n)$, with edges between $(v,n)$ and $(w,n)$ whenever $v\sim_Gw$, and between $(v,n)$ and $(v,n+1)$. Independently retain its edges with probability $p$.

Does there exist a strongly amenable $G$ and numbers $0<p_1<p_2<1$ such that, for each fixed $p\in[p_1,p_2]$, this percolation has infinitely many infinite open components almost surely? Can such an example be obtained with polynomial volume growth, meaning that for some vertex $o$ and constants $C,D<\infty$,
\[
|B_G(o,r)|\le Cr^D\qquad(r\ge1)?
\]
The question is existential over graphs, rather than an assertion that every strongly amenable graph has such an interval. The polynomial-growth part asks for an example in that narrower class.

Requiring a nondegenerate interval distinguishes the question from exceptional behavior at one parameter. Since these products need not be transitive, uniqueness of their infinite cluster need not be monotone in $p$; no equivalence with a simple $p_c<p_u$ criterion is assumed. The two product directions and the edges within each layer use the same Bernoulli parameter.
\subsection{Short English statement}
Can a product of a strongly amenable graph with the line have infinitely many infinite clusters throughout an interval of percolation parameters?
\subsection{Sources}
\begin{itemize}
\item[S1] ULAM.AI and the UnsolvedMath Contributors, supplied problems.json, record 10000042 (AMR-099-0042), AMR Open Problem Lists. Catalogue identity and source transcription; CC BY 4.0.
\url{https://huggingface.co/datasets/ulamai/UnsolvedMath}.
\item[S2] Benjamini - Coarse Geometry and Randomness (Saint-Flour notes), Open problem 9.48, PDF page 75. Original source indexed by AMR-099-0042.
\url{https://www.wisdom.weizmann.ac.il/~itai/stflouraug24.pdf#page=75}.
\item[S3] I. Benjamini, Coarse Geometry and Randomness, primary Saint-Flour notes; the numbered question and its surrounding definitions.
\url{https://arquivo.pt/noFrame/replay/20201231041548id_/http://www.wisdom.weizmann.ac.il/~itai/stflouraug24.pdf}.
\end{itemize}
\end{document}
